
\documentclass{article}
\usepackage{colortbl}
\usepackage{makecell}
\usepackage{multirow}
\usepackage{supertabular}

\begin{document}

\newcounter{utterance}

\centering \large Interaction Transcript for game `clean\_up', experiment `0\_easy\_3obj\_hu', episode 4 with Qwen3.6{-}35B{-}A3B{-}FP8{-}without{-}reasoning.
\vspace{24pt}

{ \footnotesize  \setcounter{utterance}{1}
\setlength{\tabcolsep}{0pt}
\begin{supertabular}{c@{$\;$}|p{.15\linewidth}@{}p{.15\linewidth}p{.15\linewidth}p{.15\linewidth}p{.15\linewidth}p{.15\linewidth}}
    \# & $\;$A & \multicolumn{4}{c}{Game Master} & $\;\:$B\\
    \hline

    \theutterance \stepcounter{utterance}  
    & & \multicolumn{4}{p{0.6\linewidth}}{
        \cellcolor[rgb]{0.9,0.9,0.9}{
            \makecell[{{p{\linewidth}}}]{
                \texttt{\tiny{[P1$\langle$GM]}}
                \texttt{Én vagyok a játékmestered, és egy együttműködésen alapuló játékot játszol az alábbi rács segítségével:\textbackslash{}n\textbackslash{}n```\textbackslash{}n 1234567} \\
\texttt{╔═╤═════╗} \\
\texttt{║◌│◌◌◌◌L║ 1} \\
\texttt{║◌│◌◌◌◌◌║ 2} \\
\texttt{║◌│◌◌◌◌◌║ 3} \\
\texttt{║◌│◌◌◌◌◌║ 4} \\
\texttt{╟─┼───┐C║ 5} \\
\texttt{║◌│◌◌◌├─╢ 6} \\
\texttt{║P│◌◌◌│◌║ 7} \\
\texttt{╚═╧═══╧═╝} \\
\texttt{\textbackslash{}n```\textbackslash{}n\textbackslash{}n* A felső szélén az x{-}koordináták jobbra növekednek, a jobb szélén pedig az y{-}koordináták lefelé növekednek.\textbackslash{}n* A következő tárgyak véletlenszerűen helyezkednek el a rácsodon: 'C', 'L', 'P'.\textbackslash{}n\textbackslash{}n} \\
\texttt{A másik játékos a táblának egy másik változatát látja, ahol a tárgyak különböző véletlenszerű helyeken vannak. Te nem látod az ő tábláját, és ő sem látja a tiédet.} \\
\\ 
\\ 
\texttt{**Cél:**} \\
\\ 
\texttt{Mindkét játékosnak el kell helyeznie a saját tábláján lévő tárgyakat úgy, hogy azonos tárgyak ugyanazokra a koordinátákra kerüljenek. Kommunikálnotok kell, hogy közös célt érjetek el.} \\
\\ 
\\ 
\texttt{**Szabályok:**} \\
\\ 
\texttt{* Minden körben pontosan egyet küldhetsz az alábbi két parancs közül:} \\
\texttt{1. `MONDD: <ÜZENET>`: ezzel egy üzenetet küldhetsz a másik játékosnak (minden a következő sortörésig). Én továbbítom neki.} \\
\texttt{2. `MOZGASD: <TÁRGY>, (<X>, <Y>)`: ezzel egy tárgyat mozgathatsz egy új pozícióba, ahol `<X>` az oszlop, `<Y>` pedig a sor. Én visszajelzek, hogy sikeres volt{-}e a lépés.} \\
\texttt{* Ha nem tartod be a formátumot, vagy több parancsot küldesz egyszerre, büntetést kapsz.} \\
\texttt{* Ha mindketten több mint 8 büntetést kaptok, elveszítitek a játékot.} \\
\texttt{* Elengedhetetlen, hogy egyeztessetek a célállapotról! Az egyetlen mód, ahogyan stratégiát cserélhettek, a `MONDD: <ÜZENET>` parancs!} \\
\\ 
\\ 
\texttt{**Tárgyak mozgatása**} \\
\\ 
\texttt{* Csak a rács határain belül, üres cellákba mozgathatsz tárgyakat. Az üres cellák '◌' szimbólummal vannak jelölve.} \\
\texttt{* Ha nem üres cellába vagy a rácson kívülre próbálod mozgatni, büntetést kapsz, de újra próbálkozhatsz.} \\
\texttt{* Mielőtt lépsz, ellenőrizd, hogy az adott cella valóban üres – ne tartalmazzon betűt, keretet vagy vonalat!} \\
\\ 
\\ 
\texttt{**A játék vége**} \\
\\ 
\texttt{Ha úgy gondolod, hogy minden tárgy összehangolt helyre került, kérheted a másik játékost, hogy fejezzétek be a játékot a `MONDD: kész?` parancs elküldésével. Ha a másik játékos kér erre, és te `MONDD: kész!`{-}vel válaszolsz, a játék véget ér.} \\
\\ 
\texttt{A játékot megnyeritek, ha az 12 körön belül véget ér. Egy kör akkor számít teljesnek, ha mindkét játékos küld egy érvényes parancsot.} \\
\\ 
\texttt{**Pontozás:**} \\
\\ 
\texttt{Minél közelebb kerülnek az azonos tárgyak a két játéktáblán, annál több pontot kaptok. A hibákért levonás jár. Vajon meg tudjátok dönteni a rekordot?} \\
\\ 
\\ 
\texttt{Kérlek, küldj egy üzenetet a másik játékosnak a játék elindításához!} \\
            }
        }
    }
    & \\ \\

    \theutterance \stepcounter{utterance}  
    & \multicolumn{4}{p{0.6\linewidth}}{
        \cellcolor[rgb]{0.9,0.9,0.9}{
            \makecell[{{p{\linewidth}}}]{
                \texttt{\tiny{[P1$\rangle$GM]}}
                \texttt{MONDD: Szia! Mivel a tábláinkon a tárgyak (C, L, P) véletlenszerűen helyezkednek el, és nem látjuk egymás tábláját, a legfontosabb feladatunk, hogy egyeztessünk egy közös koordináta{-}rendszerben és stratégiában.} \\
\\ 
\texttt{Javaslatom a következő lépések sorrendjére:} \\
\texttt{1. Először mozgassuk el az egyik tárgyat (pl. 'P') egy jól definiált, könnyen azonosítható pozícióra mindkét táblán (pl. bal felső sarok vagy egy központi pont). Mivel a rács szimmetrikus és a 'P' egyedi, ha mindketten ugyanarra a koordinátára mozgatjuk, az lesz a "referencia pontunk".} \\
\texttt{2. Miután a 'P' tárgyak egybeestek, használjuk ezt a pozíciót kiindulási pontként a többi tárgy ('C' és 'L') relatív elhelyezéséhez.} \\
\\ 
\texttt{Kezdjük azzal, hogy mindketten a 'P' tárgyat mozgatjuk a (1, 1) koordinátára (bal felső sarok, ha a rács 1{-}től indul).} \\
\\ 
\texttt{Kérem, erősítsd meg, hogy elfogadod ezt a stratégiát, vagy küldd a saját `MOZGASD: P, (1, 1)` parancsodat, ha már elkezdtél mozogni. Én is elküldöm a saját lépésemet.} \\
            }
        }
    }
    & & \\ \\

    \theutterance \stepcounter{utterance}  
    & & & \multicolumn{2}{p{0.3\linewidth}}{
        \cellcolor[rgb]{0.9,0.9,0.9}{
            \makecell[{{p{\linewidth}}}]{
                \texttt{\tiny{[GM$|$GM]}}
                \texttt{MONDD: Szia! Mivel a tábláinkon a tárgyak (C, L, P) véletlenszerűen helyezkednek el, és nem látjuk egymás tábláját, a legfontosabb feladatunk, hogy egyeztessünk egy közös koordináta{-}rendszerben és stratégiában.} \\
\\ 
\texttt{Javaslatom a következő lépések sorrendjére:} \\
\texttt{1. Először mozgassuk el az egyik tárgyat (pl. 'P') egy jól definiált, könnyen azonosítható pozícióra mindkét táblán (pl. bal felső sarok vagy egy központi pont). Mivel a rács szimmetrikus és a 'P' egyedi, ha mindketten ugyanarra a koordinátára mozgatjuk, az lesz a "referencia pontunk".} \\
\texttt{2. Miután a 'P' tárgyak egybeestek, használjuk ezt a pozíciót kiindulási pontként a többi tárgy ('C' és 'L') relatív elhelyezéséhez.} \\
\\ 
\texttt{Kezdjük azzal, hogy mindketten a 'P' tárgyat mozgatjuk a (1, 1) koordinátára (bal felső sarok, ha a rács 1{-}től indul).} \\
\\ 
\texttt{Kérem, erősítsd meg, hogy elfogadod ezt a stratégiát, vagy küldd a saját `MOZGASD: P, (1, 1)` parancsodat, ha már elkezdtél mozogni. Én is elküldöm a saját lépésemet.} \\
            }
        }
    }
    & & \\ \\

    \theutterance \stepcounter{utterance}  
    & & & \multicolumn{2}{p{0.3\linewidth}}{
        \cellcolor[rgb]{0.9,0.9,0.9}{
            \makecell[{{p{\linewidth}}}]{
                \texttt{\tiny{[GM$|$GM]}}
                \texttt{Invalid response format: MONDD: Szia! Mivel a tábláinkon a tárgyak (C, L, P) véletlenszerűen helyezkednek el, és nem látjuk egymás tábláját, a legfontosabb feladatunk, hogy egyeztessünk egy közös koordináta{-}rendszerben és stratégiában.} \\
\\ 
\texttt{Javaslatom a következő lépések sorrendjére:} \\
\texttt{1. Először mozgassuk el az egyik tárgyat (pl. 'P') egy jól definiált, könnyen azonosítható pozícióra mindkét táblán (pl. bal felső sarok vagy egy központi pont). Mivel a rács szimmetrikus és a 'P' egyedi, ha mindketten ugyanarra a koordinátára mozgatjuk, az lesz a "referencia pontunk".} \\
\texttt{2. Miután a 'P' tárgyak egybeestek, használjuk ezt a pozíciót kiindulási pontként a többi tárgy ('C' és 'L') relatív elhelyezéséhez.} \\
\\ 
\texttt{Kezdjük azzal, hogy mindketten a 'P' tárgyat mozgatjuk a (1, 1) koordinátára (bal felső sarok, ha a rács 1{-}től indul).} \\
\\ 
\texttt{Kérem, erősítsd meg, hogy elfogadod ezt a stratégiát, vagy küldd a saját MOZGASD: P, (1, 1) parancsodat, ha már elkezdtél mozogni. Én is elküldöm a saját lépésemet.} \\
            }
        }
    }
    & & \\ \\

    \theutterance \stepcounter{utterance}  
    & & \multicolumn{4}{p{0.6\linewidth}}{
        \cellcolor[rgb]{0.9,0.9,0.9}{
            \makecell[{{p{\linewidth}}}]{
                \texttt{\tiny{[P1$\langle$GM]}}
                \texttt{Büntetés: Az üzeneted több mint egy parancsot tartalmaz!} \\
\\ 
\texttt{Győződj meg róla, hogy a válaszod kizárólag a `MONDD: <ÜZENET>` vagy `MOZGASD: <TÁRGY>, (<X>, <Y>)` parancsot tartalmazza, és semmi mást!} \\
\\ 
\texttt{Eddig összesen 1 büntetést kaptatok a megengedett 8{-}ból/ből.} \\
\texttt{Kérlek, próbáld újra!} \\
            }
        }
    }
    & \\ \\

    \theutterance \stepcounter{utterance}  
    & \multicolumn{4}{p{0.6\linewidth}}{
        \cellcolor[rgb]{0.9,0.9,0.9}{
            \makecell[{{p{\linewidth}}}]{
                \texttt{\tiny{[P1$\rangle$GM]}}
                \texttt{MONDD: Bocsánat a hibáért. Kezdjük a stratégiát: mozgassuk mindketten a 'P' tárgyat a (1,1) pozícióba, hogy egyeztessük a koordináta{-}rendszerünket.} \\
            }
        }
    }
    & & \\ \\

    \theutterance \stepcounter{utterance}  
    & & & \multicolumn{2}{p{0.3\linewidth}}{
        \cellcolor[rgb]{0.9,0.9,0.9}{
            \makecell[{{p{\linewidth}}}]{
                \texttt{\tiny{[GM$|$GM]}}
                \texttt{MONDD: Bocsánat a hibáért. Kezdjük a stratégiát: mozgassuk mindketten a 'P' tárgyat a (1,1) pozícióba, hogy egyeztessük a koordináta{-}rendszerünket.} \\
            }
        }
    }
    & & \\ \\

    \theutterance \stepcounter{utterance}  
    & & \multicolumn{4}{p{0.6\linewidth}}{
        \cellcolor[rgb]{0.9,0.9,0.9}{
            \makecell[{{p{\linewidth}}}]{
                \texttt{\tiny{[GM$\rangle$P2]}}
                \texttt{Én vagyok a játékmestered, és egy együttműködésen alapuló játékot játszol az alábbi rács segítségével:\textbackslash{}n\textbackslash{}n```\textbackslash{}n 1234567} \\
\texttt{╔═╤═════╗} \\
\texttt{║◌│◌◌◌◌◌║ 1} \\
\texttt{║◌│◌◌PL◌║ 2} \\
\texttt{║◌│◌◌◌◌◌║ 3} \\
\texttt{║◌│◌◌◌◌◌║ 4} \\
\texttt{╟─┼───┐◌║ 5} \\
\texttt{║◌│◌◌◌├─╢ 6} \\
\texttt{║◌│◌C◌│◌║ 7} \\
\texttt{╚═╧═══╧═╝} \\
\texttt{\textbackslash{}n```\textbackslash{}n\textbackslash{}n* A felső szélén az x{-}koordináták jobbra növekednek, a jobb szélén pedig az y{-}koordináták lefelé növekednek.\textbackslash{}n* A következő tárgyak véletlenszerűen helyezkednek el a rácsodon: 'C', 'L', 'P'.\textbackslash{}n\textbackslash{}n} \\
\texttt{A másik játékos a táblának egy másik változatát látja, ahol a tárgyak különböző véletlenszerű helyeken vannak. Te nem látod az ő tábláját, és ő sem látja a tiédet.} \\
\\ 
\\ 
\texttt{**Cél:**} \\
\\ 
\texttt{Mindkét játékosnak el kell helyeznie a saját tábláján lévő tárgyakat úgy, hogy azonos tárgyak ugyanazokra a koordinátákra kerüljenek. Kommunikálnotok kell, hogy közös célt érjetek el.} \\
\\ 
\\ 
\texttt{**Szabályok:**} \\
\\ 
\texttt{* Minden körben pontosan egyet küldhetsz az alábbi két parancs közül:} \\
\texttt{1. `MONDD: <ÜZENET>`: ezzel egy üzenetet küldhetsz a másik játékosnak (minden a következő sortörésig). Én továbbítom neki.} \\
\texttt{2. `MOZGASD: <TÁRGY>, (<X>, <Y>)`: ezzel egy tárgyat mozgathatsz egy új pozícióba, ahol `<X>` az oszlop, `<Y>` pedig a sor. Én visszajelzek, hogy sikeres volt{-}e a lépés.} \\
\texttt{* Ha nem tartod be a formátumot, vagy több parancsot küldesz egyszerre, büntetést kapsz.} \\
\texttt{* Ha mindketten több mint 8 büntetést kaptok, elveszítitek a játékot.} \\
\texttt{* Elengedhetetlen, hogy egyeztessetek a célállapotról! Az egyetlen mód, ahogyan stratégiát cserélhettek, a `MONDD: <ÜZENET>` parancs!} \\
\\ 
\\ 
\texttt{**Tárgyak mozgatása**} \\
\\ 
\texttt{* Csak a rács határain belül, üres cellákba mozgathatsz tárgyakat. Az üres cellák '◌' szimbólummal vannak jelölve.} \\
\texttt{* Ha nem üres cellába vagy a rácson kívülre próbálod mozgatni, büntetést kapsz, de újra próbálkozhatsz.} \\
\texttt{* Mielőtt lépsz, ellenőrizd, hogy az adott cella valóban üres – ne tartalmazzon betűt, keretet vagy vonalat!} \\
\\ 
\\ 
\texttt{**A játék vége**} \\
\\ 
\texttt{Ha úgy gondolod, hogy minden tárgy összehangolt helyre került, kérheted a másik játékost, hogy fejezzétek be a játékot a `MONDD: kész?` parancs elküldésével. Ha a másik játékos kér erre, és te `MONDD: kész!`{-}vel válaszolsz, a játék véget ér.} \\
\\ 
\texttt{A játékot megnyeritek, ha az 12 körön belül véget ér. Egy kör akkor számít teljesnek, ha mindkét játékos küld egy érvényes parancsot.} \\
\\ 
\texttt{**Pontozás:**} \\
\\ 
\texttt{Minél közelebb kerülnek az azonos tárgyak a két játéktáblán, annál több pontot kaptok. A hibákért levonás jár. Vajon meg tudjátok dönteni a rekordot?} \\
\\ 
\texttt{A másik játékos ezzel az üzenettel indította el a játékot:} \\
\\ 
\texttt{"Bocsánat a hibáért. Kezdjük a stratégiát: mozgassuk mindketten a 'P' tárgyat a (1,1) pozícióba, hogy egyeztessük a koordináta{-}rendszerünket."} \\
\\ 
\texttt{Mi az első parancsod?} \\
            }
        }
    }
    & \\ \\

    \theutterance \stepcounter{utterance}  
    & & & \multicolumn{4}{p{0.6\linewidth}}{
        \cellcolor[rgb]{0.9,0.9,0.9}{
            \makecell[{{p{\linewidth}}}]{
                \texttt{\tiny{[GM$\langle$P2]}}
                \texttt{A játék szabályai szerint a cél az, hogy mindkét játékos ugyanarra a koordinátára helyezze az azonos tárgyakat. A másik játékos javaslatot tett egy közös koordináta{-}rendszer beállítására a 'P' tárgy (1,1) pozícióba mozgatásával.} \\
\\ 
\texttt{Azonban a te tábládon a 'P' tárgy jelenleg a **(4, 2)** pozícióban található (4. oszlop, 2. sor). Mivel a szabályok szerint minden körben pontosan egy parancsot küldhetsz, és a kommunikáció (`MONDD`) nem mozgat tárgyat, a leglogikusabb lépés az, hogy elfogadod a javaslatot, és megmozgatod a saját 'P' tárgyadat a megállapodott pozícióba.} \\
\\ 
\texttt{Ez egyben azt is jelzi, hogy érted a stratégiát és együttműködsz.} \\
\\ 
\texttt{Az első parancsom:} \\
\\ 
\texttt{`MOZGASD: P, (1, 1)`} \\
            }
        }
    }
     \\ \\

    \theutterance \stepcounter{utterance}  
    & & & \multicolumn{2}{p{0.3\linewidth}}{
        \cellcolor[rgb]{0.9,0.9,0.9}{
            \makecell[{{p{\linewidth}}}]{
                \texttt{\tiny{[GM$|$GM]}}
                \texttt{A játék szabályai szerint a cél az, hogy mindkét játékos ugyanarra a koordinátára helyezze az azonos tárgyakat. A másik játékos javaslatot tett egy közös koordináta{-}rendszer beállítására a 'P' tárgy (1,1) pozícióba mozgatásával.} \\
\\ 
\texttt{Azonban a te tábládon a 'P' tárgy jelenleg a **(4, 2)** pozícióban található (4. oszlop, 2. sor). Mivel a szabályok szerint minden körben pontosan egy parancsot küldhetsz, és a kommunikáció (`MONDD`) nem mozgat tárgyat, a leglogikusabb lépés az, hogy elfogadod a javaslatot, és megmozgatod a saját 'P' tárgyadat a megállapodott pozícióba.} \\
\\ 
\texttt{Ez egyben azt is jelzi, hogy érted a stratégiát és együttműködsz.} \\
\\ 
\texttt{Az első parancsom:} \\
\\ 
\texttt{`MOZGASD: P, (1, 1)`} \\
            }
        }
    }
    & & \\ \\

    \theutterance \stepcounter{utterance}  
    & & & \multicolumn{2}{p{0.3\linewidth}}{
        \cellcolor[rgb]{0.9,0.9,0.9}{
            \makecell[{{p{\linewidth}}}]{
                \texttt{\tiny{[GM$|$GM]}}
                \texttt{Invalid format: head is not empty: 'A játék szabályai szerint a cél az, hogy mindkét játékos ugyanarra a koordinátára helyezze az azonos tárgyakat. A másik játékos javaslatot tett egy közös koordináta{-}rendszer beállítására a 'P' tárgy (1,1) pozícióba mozgatásával.} \\
\\ 
\texttt{Azonban a te tábládon a 'P' tárgy jelenleg a **(4, 2)** pozícióban található (4. oszlop, 2. sor). Mivel a szabályok szerint minden körben pontosan egy parancsot küldhetsz, és a kommunikáció (MONDD) nem mozgat tárgyat, a leglogikusabb lépés az, hogy elfogadod a javaslatot, és megmozgatod a saját 'P' tárgyadat a megállapodott pozícióba.} \\
\\ 
\texttt{Ez egyben azt is jelzi, hogy érted a stratégiát és együttműködsz.} \\
\\ 
\texttt{Az első parancsom:} \\
\\ 
\texttt{'} \\
            }
        }
    }
    & & \\ \\

    \theutterance \stepcounter{utterance}  
    & & \multicolumn{4}{p{0.6\linewidth}}{
        \cellcolor[rgb]{0.9,0.9,0.9}{
            \makecell[{{p{\linewidth}}}]{
                \texttt{\tiny{[GM$\rangle$P2]}}
                \texttt{Büntetés: Az üzeneted nem tartalmazhat semmit a parancs előtt!} \\
\\ 
\texttt{Győződj meg róla, hogy a válaszod kizárólag a `MONDD: <ÜZENET>` vagy `MOZGASD: <TÁRGY>, (<X>, <Y>)` parancsot tartalmazza, és semmi mást!} \\
\\ 
\texttt{Eddig összesen 2 büntetést kaptatok a megengedett 8{-}ból/ből.} \\
\texttt{Kérlek, próbáld újra!} \\
            }
        }
    }
    & \\ \\

    \theutterance \stepcounter{utterance}  
    & & & \multicolumn{4}{p{0.6\linewidth}}{
        \cellcolor[rgb]{0.9,0.9,0.9}{
            \makecell[{{p{\linewidth}}}]{
                \texttt{\tiny{[GM$\langle$P2]}}
                \texttt{MOZGASD: P, (1, 1)} \\
            }
        }
    }
     \\ \\

    \theutterance \stepcounter{utterance}  
    & & & \multicolumn{2}{p{0.3\linewidth}}{
        \cellcolor[rgb]{0.9,0.9,0.9}{
            \makecell[{{p{\linewidth}}}]{
                \texttt{\tiny{[GM$|$GM]}}
                \texttt{MOZGASD: P, (1, 1)} \\
            }
        }
    }
    & & \\ \\

    \theutterance \stepcounter{utterance}  
    & & & \multicolumn{2}{p{0.3\linewidth}}{
        \cellcolor[rgb]{0.9,0.9,0.9}{
            \makecell[{{p{\linewidth}}}]{
                \texttt{\tiny{[GM$|$GM]}}
                \texttt{P moved to (1, 1)} \\
            }
        }
    }
    & & \\ \\

    \theutterance \stepcounter{utterance}  
    & & \multicolumn{4}{p{0.6\linewidth}}{
        \cellcolor[rgb]{0.9,0.9,0.9}{
            \makecell[{{p{\linewidth}}}]{
                \texttt{\tiny{[P1$\langle$GM]}}
                \texttt{Az üzeneted továbbítva lett a másik játékosnak.} \\
\\ 
\\ 
\texttt{Jelenleg a(z) 1. körnél tartotok a maximum 12 körből.} \\
\texttt{Eddig összesen 2 büntetést kaptatok a megengedett 8{-}ból/ből.} \\
\texttt{A másik játékos mozgatott egy tárgyat a saját rácsán.} \\
\\ 
\\ 
\texttt{Mi legyen a következő parancsod?} \\
            }
        }
    }
    & \\ \\

    \theutterance \stepcounter{utterance}  
    & \multicolumn{4}{p{0.6\linewidth}}{
        \cellcolor[rgb]{0.9,0.9,0.9}{
            \makecell[{{p{\linewidth}}}]{
                \texttt{\tiny{[P1$\rangle$GM]}}
                \texttt{MOZGASD: P, (1, 1)} \\
            }
        }
    }
    & & \\ \\

    \theutterance \stepcounter{utterance}  
    & & & \multicolumn{2}{p{0.3\linewidth}}{
        \cellcolor[rgb]{0.9,0.9,0.9}{
            \makecell[{{p{\linewidth}}}]{
                \texttt{\tiny{[GM$|$GM]}}
                \texttt{MOZGASD: P, (1, 1)} \\
            }
        }
    }
    & & \\ \\

    \theutterance \stepcounter{utterance}  
    & & & \multicolumn{2}{p{0.3\linewidth}}{
        \cellcolor[rgb]{0.9,0.9,0.9}{
            \makecell[{{p{\linewidth}}}]{
                \texttt{\tiny{[GM$|$GM]}}
                \texttt{P moved to (1, 1)} \\
            }
        }
    }
    & & \\ \\

    \theutterance \stepcounter{utterance}  
    & & \multicolumn{4}{p{0.6\linewidth}}{
        \cellcolor[rgb]{0.9,0.9,0.9}{
            \makecell[{{p{\linewidth}}}]{
                \texttt{\tiny{[GM$\rangle$P2]}}
                \texttt{A(z) 'P' sikeresen át lett mozgatva ide: (1,1). A frissített rácsod így néz ki:} \\
\texttt{```} \\
\texttt{ 1234567} \\
\texttt{╔═╤═════╗} \\
\texttt{║P│◌◌◌◌◌║ 1} \\
\texttt{║◌│◌◌◌L◌║ 2} \\
\texttt{║◌│◌◌◌◌◌║ 3} \\
\texttt{║◌│◌◌◌◌◌║ 4} \\
\texttt{╟─┼───┐◌║ 5} \\
\texttt{║◌│◌◌◌├─╢ 6} \\
\texttt{║◌│◌C◌│◌║ 7} \\
\texttt{╚═╧═══╧═╝} \\
\texttt{```} \\
\texttt{Jelenleg a(z) 2. körnél tartotok a maximum 12 körből.} \\
\texttt{Eddig összesen 2 büntetést kaptatok a megengedett 8{-}ból/ből.} \\
\texttt{A másik játékos mozgatott egy tárgyat a saját rácsán.} \\
\\ 
\\ 
\texttt{Mi legyen a következő parancsod?} \\
            }
        }
    }
    & \\ \\

    \theutterance \stepcounter{utterance}  
    & & & \multicolumn{4}{p{0.6\linewidth}}{
        \cellcolor[rgb]{0.9,0.9,0.9}{
            \makecell[{{p{\linewidth}}}]{
                \texttt{\tiny{[GM$\langle$P2]}}
                \texttt{MOZGASD: L, (1, 1)} \\
            }
        }
    }
     \\ \\

    \theutterance \stepcounter{utterance}  
    & & & \multicolumn{2}{p{0.3\linewidth}}{
        \cellcolor[rgb]{0.9,0.9,0.9}{
            \makecell[{{p{\linewidth}}}]{
                \texttt{\tiny{[GM$|$GM]}}
                \texttt{MOZGASD: L, (1, 1)} \\
            }
        }
    }
    & & \\ \\

    \theutterance \stepcounter{utterance}  
    & & & \multicolumn{2}{p{0.3\linewidth}}{
        \cellcolor[rgb]{0.9,0.9,0.9}{
            \makecell[{{p{\linewidth}}}]{
                \texttt{\tiny{[GM$|$GM]}}
                \texttt{Büntetés: a (1,1) pozíció nem üres, tartalmazza: 'P'.} \\
\\ 
\texttt{Eddig összesen 3 büntetést kaptatok a megengedett 8{-}ból/ből.} \\
\texttt{Kérlek, próbáld újra!} \\
            }
        }
    }
    & & \\ \\

    \theutterance \stepcounter{utterance}  
    & & \multicolumn{4}{p{0.6\linewidth}}{
        \cellcolor[rgb]{0.9,0.9,0.9}{
            \makecell[{{p{\linewidth}}}]{
                \texttt{\tiny{[GM$\rangle$P2]}}
                \texttt{Büntetés: a (1,1) pozíció nem üres, tartalmazza: 'P'.} \\
\\ 
\texttt{Eddig összesen 3 büntetést kaptatok a megengedett 8{-}ból/ből.} \\
\texttt{Kérlek, próbáld újra!} \\
            }
        }
    }
    & \\ \\

    \theutterance \stepcounter{utterance}  
    & & & \multicolumn{4}{p{0.6\linewidth}}{
        \cellcolor[rgb]{0.9,0.9,0.9}{
            \makecell[{{p{\linewidth}}}]{
                \texttt{\tiny{[GM$\langle$P2]}}
                \texttt{MOZGASD: L, (2, 1)} \\
            }
        }
    }
     \\ \\

    \theutterance \stepcounter{utterance}  
    & & & \multicolumn{2}{p{0.3\linewidth}}{
        \cellcolor[rgb]{0.9,0.9,0.9}{
            \makecell[{{p{\linewidth}}}]{
                \texttt{\tiny{[GM$|$GM]}}
                \texttt{MOZGASD: L, (2, 1)} \\
            }
        }
    }
    & & \\ \\

    \theutterance \stepcounter{utterance}  
    & & & \multicolumn{2}{p{0.3\linewidth}}{
        \cellcolor[rgb]{0.9,0.9,0.9}{
            \makecell[{{p{\linewidth}}}]{
                \texttt{\tiny{[GM$|$GM]}}
                \texttt{Büntetés: a (2,1) pozíció nem üres, tartalmazza: '│'.} \\
\\ 
\texttt{Eddig összesen 4 büntetést kaptatok a megengedett 8{-}ból/ből.} \\
\texttt{Kérlek, próbáld újra!} \\
            }
        }
    }
    & & \\ \\

    \theutterance \stepcounter{utterance}  
    & & \multicolumn{4}{p{0.6\linewidth}}{
        \cellcolor[rgb]{0.9,0.9,0.9}{
            \makecell[{{p{\linewidth}}}]{
                \texttt{\tiny{[GM$\rangle$P2]}}
                \texttt{Büntetés: a (2,1) pozíció nem üres, tartalmazza: '│'.} \\
\\ 
\texttt{Eddig összesen 4 büntetést kaptatok a megengedett 8{-}ból/ből.} \\
\texttt{Kérlek, próbáld újra!} \\
            }
        }
    }
    & \\ \\

    \theutterance \stepcounter{utterance}  
    & & & \multicolumn{4}{p{0.6\linewidth}}{
        \cellcolor[rgb]{0.9,0.9,0.9}{
            \makecell[{{p{\linewidth}}}]{
                \texttt{\tiny{[GM$\langle$P2]}}
                \texttt{MOZGASD: L, (3, 1)} \\
            }
        }
    }
     \\ \\

    \theutterance \stepcounter{utterance}  
    & & & \multicolumn{2}{p{0.3\linewidth}}{
        \cellcolor[rgb]{0.9,0.9,0.9}{
            \makecell[{{p{\linewidth}}}]{
                \texttt{\tiny{[GM$|$GM]}}
                \texttt{MOZGASD: L, (3, 1)} \\
            }
        }
    }
    & & \\ \\

    \theutterance \stepcounter{utterance}  
    & & & \multicolumn{2}{p{0.3\linewidth}}{
        \cellcolor[rgb]{0.9,0.9,0.9}{
            \makecell[{{p{\linewidth}}}]{
                \texttt{\tiny{[GM$|$GM]}}
                \texttt{L moved to (3, 1)} \\
            }
        }
    }
    & & \\ \\

    \theutterance \stepcounter{utterance}  
    & & \multicolumn{4}{p{0.6\linewidth}}{
        \cellcolor[rgb]{0.9,0.9,0.9}{
            \makecell[{{p{\linewidth}}}]{
                \texttt{\tiny{[P1$\langle$GM]}}
                \texttt{A(z) 'P' sikeresen át lett mozgatva ide: (1,1). A frissített rácsod így néz ki:} \\
\texttt{```} \\
\texttt{ 1234567} \\
\texttt{╔═╤═════╗} \\
\texttt{║P│◌◌◌◌L║ 1} \\
\texttt{║◌│◌◌◌◌◌║ 2} \\
\texttt{║◌│◌◌◌◌◌║ 3} \\
\texttt{║◌│◌◌◌◌◌║ 4} \\
\texttt{╟─┼───┐C║ 5} \\
\texttt{║◌│◌◌◌├─╢ 6} \\
\texttt{║◌│◌◌◌│◌║ 7} \\
\texttt{╚═╧═══╧═╝} \\
\texttt{```} \\
\texttt{Jelenleg a(z) 2. körnél tartotok a maximum 12 körből.} \\
\texttt{Eddig összesen 4 büntetést kaptatok a megengedett 8{-}ból/ből.} \\
\texttt{A másik játékos mozgatott egy tárgyat a saját rácsán.} \\
\\ 
\\ 
\texttt{Mi legyen a következő parancsod?} \\
            }
        }
    }
    & \\ \\

    \theutterance \stepcounter{utterance}  
    & \multicolumn{4}{p{0.6\linewidth}}{
        \cellcolor[rgb]{0.9,0.9,0.9}{
            \makecell[{{p{\linewidth}}}]{
                \texttt{\tiny{[P1$\rangle$GM]}}
                \texttt{MONDD: A 'P' tárgyam most a (1,1){-}en van. Mivel a másik játékos is mozgott, feltételezem, hogy ő is a 'P'{-}t mozgatta, és sikerrel járt. Most mozgassuk a 'L' tárgyat. A jelenlegi táblán a 'L' a (6,1){-}en van. Mozgassuk a 'L'{-}t a (2,1) pozícióba, hogy közel legyen a 'P'{-}hez, de ne ütközzön vele.} \\
            }
        }
    }
    & & \\ \\

    \theutterance \stepcounter{utterance}  
    & & & \multicolumn{2}{p{0.3\linewidth}}{
        \cellcolor[rgb]{0.9,0.9,0.9}{
            \makecell[{{p{\linewidth}}}]{
                \texttt{\tiny{[GM$|$GM]}}
                \texttt{MONDD: A 'P' tárgyam most a (1,1){-}en van. Mivel a másik játékos is mozgott, feltételezem, hogy ő is a 'P'{-}t mozgatta, és sikerrel járt. Most mozgassuk a 'L' tárgyat. A jelenlegi táblán a 'L' a (6,1){-}en van. Mozgassuk a 'L'{-}t a (2,1) pozícióba, hogy közel legyen a 'P'{-}hez, de ne ütközzön vele.} \\
            }
        }
    }
    & & \\ \\

    \theutterance \stepcounter{utterance}  
    & & \multicolumn{4}{p{0.6\linewidth}}{
        \cellcolor[rgb]{0.9,0.9,0.9}{
            \makecell[{{p{\linewidth}}}]{
                \texttt{\tiny{[GM$\rangle$P2]}}
                \texttt{A(z) 'L' sikeresen át lett mozgatva ide: (3,1). A frissített rácsod így néz ki:} \\
\texttt{```} \\
\texttt{ 1234567} \\
\texttt{╔═╤═════╗} \\
\texttt{║P│L◌◌◌◌║ 1} \\
\texttt{║◌│◌◌◌◌◌║ 2} \\
\texttt{║◌│◌◌◌◌◌║ 3} \\
\texttt{║◌│◌◌◌◌◌║ 4} \\
\texttt{╟─┼───┐◌║ 5} \\
\texttt{║◌│◌◌◌├─╢ 6} \\
\texttt{║◌│◌C◌│◌║ 7} \\
\texttt{╚═╧═══╧═╝} \\
\texttt{```} \\
\texttt{Jelenleg a(z) 3. körnél tartotok a maximum 12 körből.} \\
\texttt{Eddig összesen 4 büntetést kaptatok a megengedett 8{-}ból/ből.} \\
\texttt{A másik játékos ezt az üzenetet küldte:} \\
\\ 
\\ 
\texttt{"A 'P' tárgyam most a (1,1){-}en van. Mivel a másik játékos is mozgott, feltételezem, hogy ő is a 'P'{-}t mozgatta, és sikerrel járt. Most mozgassuk a 'L' tárgyat. A jelenlegi táblán a 'L' a (6,1){-}en van. Mozgassuk a 'L'{-}t a (2,1) pozícióba, hogy közel legyen a 'P'{-}hez, de ne ütközzön vele."} \\
\\ 
\\ 
\texttt{Mi legyen a következő parancsod?} \\
            }
        }
    }
    & \\ \\

    \theutterance \stepcounter{utterance}  
    & & & \multicolumn{4}{p{0.6\linewidth}}{
        \cellcolor[rgb]{0.9,0.9,0.9}{
            \makecell[{{p{\linewidth}}}]{
                \texttt{\tiny{[GM$\langle$P2]}}
                \texttt{MOZGASD: C, (1, 1)} \\
            }
        }
    }
     \\ \\

    \theutterance \stepcounter{utterance}  
    & & & \multicolumn{2}{p{0.3\linewidth}}{
        \cellcolor[rgb]{0.9,0.9,0.9}{
            \makecell[{{p{\linewidth}}}]{
                \texttt{\tiny{[GM$|$GM]}}
                \texttt{MOZGASD: C, (1, 1)} \\
            }
        }
    }
    & & \\ \\

    \theutterance \stepcounter{utterance}  
    & & & \multicolumn{2}{p{0.3\linewidth}}{
        \cellcolor[rgb]{0.9,0.9,0.9}{
            \makecell[{{p{\linewidth}}}]{
                \texttt{\tiny{[GM$|$GM]}}
                \texttt{Büntetés: a (1,1) pozíció nem üres, tartalmazza: 'P'.} \\
\\ 
\texttt{Eddig összesen 5 büntetést kaptatok a megengedett 8{-}ból/ből.} \\
\texttt{Kérlek, próbáld újra!} \\
            }
        }
    }
    & & \\ \\

    \theutterance \stepcounter{utterance}  
    & & \multicolumn{4}{p{0.6\linewidth}}{
        \cellcolor[rgb]{0.9,0.9,0.9}{
            \makecell[{{p{\linewidth}}}]{
                \texttt{\tiny{[GM$\rangle$P2]}}
                \texttt{Büntetés: a (1,1) pozíció nem üres, tartalmazza: 'P'.} \\
\\ 
\texttt{Eddig összesen 5 büntetést kaptatok a megengedett 8{-}ból/ből.} \\
\texttt{Kérlek, próbáld újra!} \\
            }
        }
    }
    & \\ \\

    \theutterance \stepcounter{utterance}  
    & & & \multicolumn{4}{p{0.6\linewidth}}{
        \cellcolor[rgb]{0.9,0.9,0.9}{
            \makecell[{{p{\linewidth}}}]{
                \texttt{\tiny{[GM$\langle$P2]}}
                \texttt{MOZGASD: C, (2, 1)} \\
            }
        }
    }
     \\ \\

    \theutterance \stepcounter{utterance}  
    & & & \multicolumn{2}{p{0.3\linewidth}}{
        \cellcolor[rgb]{0.9,0.9,0.9}{
            \makecell[{{p{\linewidth}}}]{
                \texttt{\tiny{[GM$|$GM]}}
                \texttt{MOZGASD: C, (2, 1)} \\
            }
        }
    }
    & & \\ \\

    \theutterance \stepcounter{utterance}  
    & & & \multicolumn{2}{p{0.3\linewidth}}{
        \cellcolor[rgb]{0.9,0.9,0.9}{
            \makecell[{{p{\linewidth}}}]{
                \texttt{\tiny{[GM$|$GM]}}
                \texttt{Büntetés: a (2,1) pozíció nem üres, tartalmazza: '│'.} \\
\\ 
\texttt{Eddig összesen 6 büntetést kaptatok a megengedett 8{-}ból/ből.} \\
\texttt{Kérlek, próbáld újra!} \\
            }
        }
    }
    & & \\ \\

    \theutterance \stepcounter{utterance}  
    & & \multicolumn{4}{p{0.6\linewidth}}{
        \cellcolor[rgb]{0.9,0.9,0.9}{
            \makecell[{{p{\linewidth}}}]{
                \texttt{\tiny{[GM$\rangle$P2]}}
                \texttt{Büntetés: a (2,1) pozíció nem üres, tartalmazza: '│'.} \\
\\ 
\texttt{Eddig összesen 6 büntetést kaptatok a megengedett 8{-}ból/ből.} \\
\texttt{Kérlek, próbáld újra!} \\
            }
        }
    }
    & \\ \\

    \theutterance \stepcounter{utterance}  
    & & & \multicolumn{4}{p{0.6\linewidth}}{
        \cellcolor[rgb]{0.9,0.9,0.9}{
            \makecell[{{p{\linewidth}}}]{
                \texttt{\tiny{[GM$\langle$P2]}}
                \texttt{MOZGASD: C, (3, 1)} \\
            }
        }
    }
     \\ \\

    \theutterance \stepcounter{utterance}  
    & & & \multicolumn{2}{p{0.3\linewidth}}{
        \cellcolor[rgb]{0.9,0.9,0.9}{
            \makecell[{{p{\linewidth}}}]{
                \texttt{\tiny{[GM$|$GM]}}
                \texttt{MOZGASD: C, (3, 1)} \\
            }
        }
    }
    & & \\ \\

    \theutterance \stepcounter{utterance}  
    & & & \multicolumn{2}{p{0.3\linewidth}}{
        \cellcolor[rgb]{0.9,0.9,0.9}{
            \makecell[{{p{\linewidth}}}]{
                \texttt{\tiny{[GM$|$GM]}}
                \texttt{Büntetés: a (3,1) pozíció nem üres, tartalmazza: 'L'.} \\
\\ 
\texttt{Eddig összesen 7 büntetést kaptatok a megengedett 8{-}ból/ből.} \\
\texttt{Kérlek, próbáld újra!} \\
            }
        }
    }
    & & \\ \\

    \theutterance \stepcounter{utterance}  
    & & \multicolumn{4}{p{0.6\linewidth}}{
        \cellcolor[rgb]{0.9,0.9,0.9}{
            \makecell[{{p{\linewidth}}}]{
                \texttt{\tiny{[GM$\rangle$P2]}}
                \texttt{Büntetés: a (3,1) pozíció nem üres, tartalmazza: 'L'.} \\
\\ 
\texttt{Eddig összesen 7 büntetést kaptatok a megengedett 8{-}ból/ből.} \\
\texttt{Kérlek, próbáld újra!} \\
            }
        }
    }
    & \\ \\

    \theutterance \stepcounter{utterance}  
    & & & \multicolumn{4}{p{0.6\linewidth}}{
        \cellcolor[rgb]{0.9,0.9,0.9}{
            \makecell[{{p{\linewidth}}}]{
                \texttt{\tiny{[GM$\langle$P2]}}
                \texttt{MOZGASD: C, (4, 1)} \\
            }
        }
    }
     \\ \\

    \theutterance \stepcounter{utterance}  
    & & & \multicolumn{2}{p{0.3\linewidth}}{
        \cellcolor[rgb]{0.9,0.9,0.9}{
            \makecell[{{p{\linewidth}}}]{
                \texttt{\tiny{[GM$|$GM]}}
                \texttt{MOZGASD: C, (4, 1)} \\
            }
        }
    }
    & & \\ \\

    \theutterance \stepcounter{utterance}  
    & & & \multicolumn{2}{p{0.3\linewidth}}{
        \cellcolor[rgb]{0.9,0.9,0.9}{
            \makecell[{{p{\linewidth}}}]{
                \texttt{\tiny{[GM$|$GM]}}
                \texttt{C moved to (4, 1)} \\
            }
        }
    }
    & & \\ \\

    \theutterance \stepcounter{utterance}  
    & & \multicolumn{4}{p{0.6\linewidth}}{
        \cellcolor[rgb]{0.9,0.9,0.9}{
            \makecell[{{p{\linewidth}}}]{
                \texttt{\tiny{[P1$\langle$GM]}}
                \texttt{Az üzeneted továbbítva lett a másik játékosnak.} \\
\\ 
\\ 
\texttt{Jelenleg a(z) 3. körnél tartotok a maximum 12 körből.} \\
\texttt{Eddig összesen 7 büntetést kaptatok a megengedett 8{-}ból/ből.} \\
\texttt{A másik játékos mozgatott egy tárgyat a saját rácsán.} \\
\\ 
\\ 
\texttt{Mi legyen a következő parancsod?} \\
            }
        }
    }
    & \\ \\

    \theutterance \stepcounter{utterance}  
    & \multicolumn{4}{p{0.6\linewidth}}{
        \cellcolor[rgb]{0.9,0.9,0.9}{
            \makecell[{{p{\linewidth}}}]{
                \texttt{\tiny{[P1$\rangle$GM]}}
                \texttt{MOZGASD: L, (2, 1)} \\
            }
        }
    }
    & & \\ \\

    \theutterance \stepcounter{utterance}  
    & & & \multicolumn{2}{p{0.3\linewidth}}{
        \cellcolor[rgb]{0.9,0.9,0.9}{
            \makecell[{{p{\linewidth}}}]{
                \texttt{\tiny{[GM$|$GM]}}
                \texttt{MOZGASD: L, (2, 1)} \\
            }
        }
    }
    & & \\ \\

    \theutterance \stepcounter{utterance}  
    & & & \multicolumn{2}{p{0.3\linewidth}}{
        \cellcolor[rgb]{0.9,0.9,0.9}{
            \makecell[{{p{\linewidth}}}]{
                \texttt{\tiny{[GM$|$GM]}}
                \texttt{Büntetés: a (2,1) pozíció nem üres, tartalmazza: '│'.} \\
\\ 
\texttt{Eddig összesen 8 büntetést kaptatok a megengedett 8{-}ból/ből.} \\
\texttt{Kérlek, próbáld újra!} \\
            }
        }
    }
    & & \\ \\

    \theutterance \stepcounter{utterance}  
    & & \multicolumn{4}{p{0.6\linewidth}}{
        \cellcolor[rgb]{0.9,0.9,0.9}{
            \makecell[{{p{\linewidth}}}]{
                \texttt{\tiny{[P1$\langle$GM]}}
                \texttt{Büntetés: a (2,1) pozíció nem üres, tartalmazza: '│'.} \\
\\ 
\texttt{Eddig összesen 8 büntetést kaptatok a megengedett 8{-}ból/ből.} \\
\texttt{Kérlek, próbáld újra!} \\
            }
        }
    }
    & \\ \\

    \theutterance \stepcounter{utterance}  
    & \multicolumn{4}{p{0.6\linewidth}}{
        \cellcolor[rgb]{0.9,0.9,0.9}{
            \makecell[{{p{\linewidth}}}]{
                \texttt{\tiny{[P1$\rangle$GM]}}
                \texttt{MOZGASD: L, (2, 2)} \\
            }
        }
    }
    & & \\ \\

    \theutterance \stepcounter{utterance}  
    & & & \multicolumn{2}{p{0.3\linewidth}}{
        \cellcolor[rgb]{0.9,0.9,0.9}{
            \makecell[{{p{\linewidth}}}]{
                \texttt{\tiny{[GM$|$GM]}}
                \texttt{MOZGASD: L, (2, 2)} \\
            }
        }
    }
    & & \\ \\

    \theutterance \stepcounter{utterance}  
    & & & \multicolumn{2}{p{0.3\linewidth}}{
        \cellcolor[rgb]{0.9,0.9,0.9}{
            \makecell[{{p{\linewidth}}}]{
                \texttt{\tiny{[GM$|$GM]}}
                \texttt{Büntetés: a (2,2) pozíció nem üres, tartalmazza: '│'.} \\
\\ 
\texttt{Eddig összesen 9 büntetést kaptatok a megengedett 8{-}ból/ből.} \\
\texttt{Kérlek, próbáld újra!} \\
            }
        }
    }
    & & \\ \\

    \theutterance \stepcounter{utterance}  
    & & \multicolumn{4}{p{0.6\linewidth}}{
        \cellcolor[rgb]{0.9,0.9,0.9}{
            \makecell[{{p{\linewidth}}}]{
                \texttt{\tiny{[P1$\langle$GM]}}
                \texttt{Büntetés: a (2,2) pozíció nem üres, tartalmazza: '│'.} \\
\\ 
\texttt{Eddig összesen 9 büntetést kaptatok a megengedett 8{-}ból/ből.} \\
\texttt{Kérlek, próbáld újra!} \\
            }
        }
    }
    & \\ \\

    \theutterance \stepcounter{utterance}  
    & \multicolumn{4}{p{0.6\linewidth}}{
        \cellcolor[rgb]{0.9,0.9,0.9}{
            \makecell[{{p{\linewidth}}}]{
                \texttt{\tiny{[P1$\rangle$GM]}}
                \texttt{MOZGASD: L, (3, 1)} \\
            }
        }
    }
    & & \\ \\

    \theutterance \stepcounter{utterance}  
    & & & \multicolumn{2}{p{0.3\linewidth}}{
        \cellcolor[rgb]{0.9,0.9,0.9}{
            \makecell[{{p{\linewidth}}}]{
                \texttt{\tiny{[GM$|$GM]}}
                \texttt{MOZGASD: L, (3, 1)} \\
            }
        }
    }
    & & \\ \\

    \theutterance \stepcounter{utterance}  
    & & & \multicolumn{2}{p{0.3\linewidth}}{
        \cellcolor[rgb]{0.9,0.9,0.9}{
            \makecell[{{p{\linewidth}}}]{
                \texttt{\tiny{[GM$|$GM]}}
                \texttt{L moved to (3, 1)} \\
            }
        }
    }
    & & \\ \\

    \theutterance \stepcounter{utterance}  
    & & \multicolumn{4}{p{0.6\linewidth}}{
        \cellcolor[rgb]{0.9,0.9,0.9}{
            \makecell[{{p{\linewidth}}}]{
                \texttt{\tiny{[GM$\rangle$P2]}}
                \texttt{A(z) 'C' sikeresen át lett mozgatva ide: (4,1). A frissített rácsod így néz ki:} \\
\texttt{```} \\
\texttt{ 1234567} \\
\texttt{╔═╤═════╗} \\
\texttt{║P│LC◌◌◌║ 1} \\
\texttt{║◌│◌◌◌◌◌║ 2} \\
\texttt{║◌│◌◌◌◌◌║ 3} \\
\texttt{║◌│◌◌◌◌◌║ 4} \\
\texttt{╟─┼───┐◌║ 5} \\
\texttt{║◌│◌◌◌├─╢ 6} \\
\texttt{║◌│◌◌◌│◌║ 7} \\
\texttt{╚═╧═══╧═╝} \\
\texttt{```} \\
\texttt{Jelenleg a(z) 4. körnél tartotok a maximum 12 körből.} \\
\texttt{Eddig összesen 9 büntetést kaptatok a megengedett 8{-}ból/ből.} \\
\texttt{A másik játékos mozgatott egy tárgyat a saját rácsán.} \\
\\ 
\\ 
\texttt{Mi legyen a következő parancsod?} \\
            }
        }
    }
    & \\ \\

    \theutterance \stepcounter{utterance}  
    & & & \multicolumn{4}{p{0.6\linewidth}}{
        \cellcolor[rgb]{0.9,0.9,0.9}{
            \makecell[{{p{\linewidth}}}]{
                \texttt{\tiny{[GM$\langle$P2]}}
                \texttt{MOZGASD: C, (5, 1)} \\
            }
        }
    }
     \\ \\

    \theutterance \stepcounter{utterance}  
    & & & \multicolumn{2}{p{0.3\linewidth}}{
        \cellcolor[rgb]{0.9,0.9,0.9}{
            \makecell[{{p{\linewidth}}}]{
                \texttt{\tiny{[GM$|$GM]}}
                \texttt{MOZGASD: C, (5, 1)} \\
            }
        }
    }
    & & \\ \\

    \theutterance \stepcounter{utterance}  
    & & & \multicolumn{2}{p{0.3\linewidth}}{
        \cellcolor[rgb]{0.9,0.9,0.9}{
            \makecell[{{p{\linewidth}}}]{
                \texttt{\tiny{[GM$|$GM]}}
                \texttt{C moved to (5, 1)} \\
            }
        }
    }
    & & \\ \\

    \theutterance \stepcounter{utterance}  
    & & \multicolumn{4}{p{0.6\linewidth}}{
        \cellcolor[rgb]{0.9,0.9,0.9}{
            \makecell[{{p{\linewidth}}}]{
                \texttt{\tiny{[P1$\langle$GM]}}
                \texttt{A(z) 'L' sikeresen át lett mozgatva ide: (3,1). A frissített rácsod így néz ki:} \\
\texttt{```} \\
\texttt{ 1234567} \\
\texttt{╔═╤═════╗} \\
\texttt{║P│L◌◌◌◌║ 1} \\
\texttt{║◌│◌◌◌◌◌║ 2} \\
\texttt{║◌│◌◌◌◌◌║ 3} \\
\texttt{║◌│◌◌◌◌◌║ 4} \\
\texttt{╟─┼───┐C║ 5} \\
\texttt{║◌│◌◌◌├─╢ 6} \\
\texttt{║◌│◌◌◌│◌║ 7} \\
\texttt{╚═╧═══╧═╝} \\
\texttt{```} \\
\texttt{Jelenleg a(z) 4. körnél tartotok a maximum 12 körből.} \\
\texttt{Eddig összesen 9 büntetést kaptatok a megengedett 8{-}ból/ből.} \\
\texttt{A másik játékos mozgatott egy tárgyat a saját rácsán.} \\
\\ 
\\ 
\texttt{Mi legyen a következő parancsod?} \\
            }
        }
    }
    & \\ \\

    \theutterance \stepcounter{utterance}  
    & \multicolumn{4}{p{0.6\linewidth}}{
        \cellcolor[rgb]{0.9,0.9,0.9}{
            \makecell[{{p{\linewidth}}}]{
                \texttt{\tiny{[P1$\rangle$GM]}}
                \texttt{MONDD: A 'L' most a (3,1){-}en van. A 'C' tárgyat még nem mozgattam. A jelenlegi táblán a 'C' a (6,5){-}ön van. Mozgassuk a 'C'{-}t a (4,1) pozícióba, hogy a sorban legyen a 'P' és 'L' mellett.} \\
            }
        }
    }
    & & \\ \\

    \theutterance \stepcounter{utterance}  
    & & & \multicolumn{2}{p{0.3\linewidth}}{
        \cellcolor[rgb]{0.9,0.9,0.9}{
            \makecell[{{p{\linewidth}}}]{
                \texttt{\tiny{[GM$|$GM]}}
                \texttt{MONDD: A 'L' most a (3,1){-}en van. A 'C' tárgyat még nem mozgattam. A jelenlegi táblán a 'C' a (6,5){-}ön van. Mozgassuk a 'C'{-}t a (4,1) pozícióba, hogy a sorban legyen a 'P' és 'L' mellett.} \\
            }
        }
    }
    & & \\ \\

    \theutterance \stepcounter{utterance}  
    & & \multicolumn{4}{p{0.6\linewidth}}{
        \cellcolor[rgb]{0.9,0.9,0.9}{
            \makecell[{{p{\linewidth}}}]{
                \texttt{\tiny{[GM$\rangle$P2]}}
                \texttt{A(z) 'C' sikeresen át lett mozgatva ide: (5,1). A frissített rácsod így néz ki:} \\
\texttt{```} \\
\texttt{ 1234567} \\
\texttt{╔═╤═════╗} \\
\texttt{║P│L◌C◌◌║ 1} \\
\texttt{║◌│◌◌◌◌◌║ 2} \\
\texttt{║◌│◌◌◌◌◌║ 3} \\
\texttt{║◌│◌◌◌◌◌║ 4} \\
\texttt{╟─┼───┐◌║ 5} \\
\texttt{║◌│◌◌◌├─╢ 6} \\
\texttt{║◌│◌◌◌│◌║ 7} \\
\texttt{╚═╧═══╧═╝} \\
\texttt{```} \\
\texttt{Jelenleg a(z) 5. körnél tartotok a maximum 12 körből.} \\
\texttt{Eddig összesen 9 büntetést kaptatok a megengedett 8{-}ból/ből.} \\
\texttt{A másik játékos ezt az üzenetet küldte:} \\
\\ 
\\ 
\texttt{"A 'L' most a (3,1){-}en van. A 'C' tárgyat még nem mozgattam. A jelenlegi táblán a 'C' a (6,5){-}ön van. Mozgassuk a 'C'{-}t a (4,1) pozícióba, hogy a sorban legyen a 'P' és 'L' mellett."} \\
\\ 
\\ 
\texttt{Mi legyen a következő parancsod?} \\
            }
        }
    }
    & \\ \\

    \theutterance \stepcounter{utterance}  
    & & & \multicolumn{4}{p{0.6\linewidth}}{
        \cellcolor[rgb]{0.9,0.9,0.9}{
            \makecell[{{p{\linewidth}}}]{
                \texttt{\tiny{[GM$\langle$P2]}}
                \texttt{MONDD: A 'C' már a (5,1){-}en van, nem a (4,1){-}en. A 'P' és 'L' is az első sorban van. Kérlek, mozgassd a 'C'{-}t a (6,1){-}re, hogy a sorban legyen, vagy egyeztessünk egy másik stratégiát.} \\
            }
        }
    }
     \\ \\

    \theutterance \stepcounter{utterance}  
    & & & \multicolumn{2}{p{0.3\linewidth}}{
        \cellcolor[rgb]{0.9,0.9,0.9}{
            \makecell[{{p{\linewidth}}}]{
                \texttt{\tiny{[GM$|$GM]}}
                \texttt{MONDD: A 'C' már a (5,1){-}en van, nem a (4,1){-}en. A 'P' és 'L' is az első sorban van. Kérlek, mozgassd a 'C'{-}t a (6,1){-}re, hogy a sorban legyen, vagy egyeztessünk egy másik stratégiát.} \\
            }
        }
    }
    & & \\ \\

    \theutterance \stepcounter{utterance}  
    & & \multicolumn{4}{p{0.6\linewidth}}{
        \cellcolor[rgb]{0.9,0.9,0.9}{
            \makecell[{{p{\linewidth}}}]{
                \texttt{\tiny{[P1$\langle$GM]}}
                \texttt{Az üzeneted továbbítva lett a másik játékosnak.} \\
\\ 
\\ 
\texttt{Jelenleg a(z) 5. körnél tartotok a maximum 12 körből.} \\
\texttt{Eddig összesen 9 büntetést kaptatok a megengedett 8{-}ból/ből.} \\
\texttt{A másik játékos ezt az üzenetet küldte:} \\
\\ 
\\ 
\texttt{"A 'C' már a (5,1){-}en van, nem a (4,1){-}en. A 'P' és 'L' is az első sorban van. Kérlek, mozgassd a 'C'{-}t a (6,1){-}re, hogy a sorban legyen, vagy egyeztessünk egy másik stratégiát."} \\
\\ 
\\ 
\texttt{Mi legyen a következő parancsod?} \\
            }
        }
    }
    & \\ \\

    \theutterance \stepcounter{utterance}  
    & \multicolumn{4}{p{0.6\linewidth}}{
        \cellcolor[rgb]{0.9,0.9,0.9}{
            \makecell[{{p{\linewidth}}}]{
                \texttt{\tiny{[P1$\rangle$GM]}}
                \texttt{MOZGASD: C, (6, 1)} \\
            }
        }
    }
    & & \\ \\

    \theutterance \stepcounter{utterance}  
    & & & \multicolumn{2}{p{0.3\linewidth}}{
        \cellcolor[rgb]{0.9,0.9,0.9}{
            \makecell[{{p{\linewidth}}}]{
                \texttt{\tiny{[GM$|$GM]}}
                \texttt{MOZGASD: C, (6, 1)} \\
            }
        }
    }
    & & \\ \\

    \theutterance \stepcounter{utterance}  
    & & & \multicolumn{2}{p{0.3\linewidth}}{
        \cellcolor[rgb]{0.9,0.9,0.9}{
            \makecell[{{p{\linewidth}}}]{
                \texttt{\tiny{[GM$|$GM]}}
                \texttt{C moved to (6, 1)} \\
            }
        }
    }
    & & \\ \\

    \theutterance \stepcounter{utterance}  
    & & \multicolumn{4}{p{0.6\linewidth}}{
        \cellcolor[rgb]{0.9,0.9,0.9}{
            \makecell[{{p{\linewidth}}}]{
                \texttt{\tiny{[GM$\rangle$P2]}}
                \texttt{Az üzeneted továbbítva lett a másik játékosnak.} \\
\\ 
\\ 
\texttt{Jelenleg a(z) 6. körnél tartotok a maximum 12 körből.} \\
\texttt{Eddig összesen 9 büntetést kaptatok a megengedett 8{-}ból/ből.} \\
\texttt{A másik játékos mozgatott egy tárgyat a saját rácsán.} \\
\\ 
\\ 
\texttt{Mi legyen a következő parancsod?} \\
            }
        }
    }
    & \\ \\

    \theutterance \stepcounter{utterance}  
    & & & \multicolumn{4}{p{0.6\linewidth}}{
        \cellcolor[rgb]{0.9,0.9,0.9}{
            \makecell[{{p{\linewidth}}}]{
                \texttt{\tiny{[GM$\langle$P2]}}
                \texttt{MOZGASD: C, (6, 1)} \\
            }
        }
    }
     \\ \\

    \theutterance \stepcounter{utterance}  
    & & & \multicolumn{2}{p{0.3\linewidth}}{
        \cellcolor[rgb]{0.9,0.9,0.9}{
            \makecell[{{p{\linewidth}}}]{
                \texttt{\tiny{[GM$|$GM]}}
                \texttt{MOZGASD: C, (6, 1)} \\
            }
        }
    }
    & & \\ \\

    \theutterance \stepcounter{utterance}  
    & & & \multicolumn{2}{p{0.3\linewidth}}{
        \cellcolor[rgb]{0.9,0.9,0.9}{
            \makecell[{{p{\linewidth}}}]{
                \texttt{\tiny{[GM$|$GM]}}
                \texttt{C moved to (6, 1)} \\
            }
        }
    }
    & & \\ \\

    \theutterance \stepcounter{utterance}  
    & & \multicolumn{4}{p{0.6\linewidth}}{
        \cellcolor[rgb]{0.9,0.9,0.9}{
            \makecell[{{p{\linewidth}}}]{
                \texttt{\tiny{[P1$\langle$GM]}}
                \texttt{A(z) 'C' sikeresen át lett mozgatva ide: (6,1). A frissített rácsod így néz ki:} \\
\texttt{```} \\
\texttt{ 1234567} \\
\texttt{╔═╤═════╗} \\
\texttt{║P│L◌◌C◌║ 1} \\
\texttt{║◌│◌◌◌◌◌║ 2} \\
\texttt{║◌│◌◌◌◌◌║ 3} \\
\texttt{║◌│◌◌◌◌◌║ 4} \\
\texttt{╟─┼───┐◌║ 5} \\
\texttt{║◌│◌◌◌├─╢ 6} \\
\texttt{║◌│◌◌◌│◌║ 7} \\
\texttt{╚═╧═══╧═╝} \\
\texttt{```} \\
\texttt{Jelenleg a(z) 6. körnél tartotok a maximum 12 körből.} \\
\texttt{Eddig összesen 9 büntetést kaptatok a megengedett 8{-}ból/ből.} \\
\texttt{A másik játékos mozgatott egy tárgyat a saját rácsán.} \\
\\ 
\\ 
\texttt{Mi legyen a következő parancsod?} \\
            }
        }
    }
    & \\ \\

    \theutterance \stepcounter{utterance}  
    & \multicolumn{4}{p{0.6\linewidth}}{
        \cellcolor[rgb]{0.9,0.9,0.9}{
            \makecell[{{p{\linewidth}}}]{
                \texttt{\tiny{[P1$\rangle$GM]}}
                \texttt{MONDD: Kiváló! A 'P', 'L' és 'C' tárgyak most mind az első sorban vannak: (1,1), (3,1) és (6,1). Mivel mindkét játékos ugyanerre a konfigurációra törekszik, és a tárgyak egyedi elhelyezkedése stabilizálódott, úgy tűnik, hogy a célállapotot elértük. Kérlek, erősítsd meg, hogy a te tábládon is ugyanerre a pozíciókra kerültek a tárgyak, és ha igen, küldd a `MONDD: kész!` üzenetet.} \\
            }
        }
    }
    & & \\ \\

    \theutterance \stepcounter{utterance}  
    & & & \multicolumn{2}{p{0.3\linewidth}}{
        \cellcolor[rgb]{0.9,0.9,0.9}{
            \makecell[{{p{\linewidth}}}]{
                \texttt{\tiny{[GM$|$GM]}}
                \texttt{MONDD: Kiváló! A 'P', 'L' és 'C' tárgyak most mind az első sorban vannak: (1,1), (3,1) és (6,1). Mivel mindkét játékos ugyanerre a konfigurációra törekszik, és a tárgyak egyedi elhelyezkedése stabilizálódott, úgy tűnik, hogy a célállapotot elértük. Kérlek, erősítsd meg, hogy a te tábládon is ugyanerre a pozíciókra kerültek a tárgyak, és ha igen, küldd a `MONDD: kész!` üzenetet.} \\
            }
        }
    }
    & & \\ \\

    \theutterance \stepcounter{utterance}  
    & & & \multicolumn{2}{p{0.3\linewidth}}{
        \cellcolor[rgb]{0.9,0.9,0.9}{
            \makecell[{{p{\linewidth}}}]{
                \texttt{\tiny{[GM$|$GM]}}
                \texttt{Invalid format: head is not empty: 'MONDD: Kiváló! A 'P', 'L' és 'C' tárgyak most mind az első sorban vannak: (1,1), (3,1) és (6,1). Mivel mindkét játékos ugyanerre a konfigurációra törekszik, és a tárgyak egyedi elhelyezkedése stabilizálódott, úgy tűnik, hogy a célállapotot elértük. Kérlek, erősítsd meg, hogy a te tábládon is ugyanerre a pozíciókra kerültek a tárgyak, és ha igen, küldd a '} \\
            }
        }
    }
    & & \\ \\

    \theutterance \stepcounter{utterance}  
    & & & \multicolumn{2}{p{0.3\linewidth}}{
        \cellcolor[rgb]{0.9,0.9,0.9}{
            \makecell[{{p{\linewidth}}}]{
                \texttt{\tiny{[GM$|$GM]}}
                \texttt{Maximum number of penalties exceeded} \\
            }
        }
    }
    & & \\ \\

    \theutterance \stepcounter{utterance}  
    & & & \multicolumn{2}{p{0.3\linewidth}}{
        \cellcolor[rgb]{0.9,0.9,0.9}{
            \makecell[{{p{\linewidth}}}]{
                \texttt{\tiny{[GM$|$GM]}}
                \texttt{Initial states:} \\
\texttt{```} \\
\texttt{Player 1:} \\
\texttt{ 1234567} \\
\texttt{╔═╤═════╗} \\
\texttt{║◌│◌◌◌◌L║ 1} \\
\texttt{║◌│◌◌◌◌◌║ 2} \\
\texttt{║◌│◌◌◌◌◌║ 3} \\
\texttt{║◌│◌◌◌◌◌║ 4} \\
\texttt{╟─┼───┐C║ 5} \\
\texttt{║◌│◌◌◌├─╢ 6} \\
\texttt{║P│◌◌◌│◌║ 7} \\
\texttt{╚═╧═══╧═╝} \\
\\ 
\\ 
\texttt{Player 2:} \\
\texttt{ 1234567} \\
\texttt{╔═╤═════╗} \\
\texttt{║◌│◌◌◌◌◌║ 1} \\
\texttt{║◌│◌◌PL◌║ 2} \\
\texttt{║◌│◌◌◌◌◌║ 3} \\
\texttt{║◌│◌◌◌◌◌║ 4} \\
\texttt{╟─┼───┐◌║ 5} \\
\texttt{║◌│◌◌◌├─╢ 6} \\
\texttt{║◌│◌C◌│◌║ 7} \\
\texttt{╚═╧═══╧═╝} \\
\\ 
\texttt{```} \\
            }
        }
    }
    & & \\ \\

    \theutterance \stepcounter{utterance}  
    & & & \multicolumn{2}{p{0.3\linewidth}}{
        \cellcolor[rgb]{0.9,0.9,0.9}{
            \makecell[{{p{\linewidth}}}]{
                \texttt{\tiny{[GM$|$GM]}}
                \texttt{End states:} \\
\texttt{```} \\
\texttt{Player 1:} \\
\texttt{ 1234567} \\
\texttt{╔═╤═════╗} \\
\texttt{║P│L◌◌C◌║ 1} \\
\texttt{║◌│◌◌◌◌◌║ 2} \\
\texttt{║◌│◌◌◌◌◌║ 3} \\
\texttt{║◌│◌◌◌◌◌║ 4} \\
\texttt{╟─┼───┐◌║ 5} \\
\texttt{║◌│◌◌◌├─╢ 6} \\
\texttt{║◌│◌◌◌│◌║ 7} \\
\texttt{╚═╧═══╧═╝} \\
\\ 
\\ 
\texttt{Player 2:} \\
\texttt{ 1234567} \\
\texttt{╔═╤═════╗} \\
\texttt{║P│L◌◌C◌║ 1} \\
\texttt{║◌│◌◌◌◌◌║ 2} \\
\texttt{║◌│◌◌◌◌◌║ 3} \\
\texttt{║◌│◌◌◌◌◌║ 4} \\
\texttt{╟─┼───┐◌║ 5} \\
\texttt{║◌│◌◌◌├─╢ 6} \\
\texttt{║◌│◌◌◌│◌║ 7} \\
\texttt{╚═╧═══╧═╝} \\
\\ 
\texttt{```} \\
            }
        }
    }
    & & \\ \\

    \theutterance \stepcounter{utterance}  
    & & & \multicolumn{2}{p{0.3\linewidth}}{
        \cellcolor[rgb]{0.9,0.9,0.9}{
            \makecell[{{p{\linewidth}}}]{
                \texttt{\tiny{[GM$|$GM]}}
                \texttt{* success: False} \\
\texttt{* lose: False} \\
\texttt{* aborted: True} \\
\texttt{{-}{-}{-}{-}{-}{-}{-}} \\
\texttt{* Shifts: 4.00} \\
\texttt{* Max Shifts: 4.00} \\
\texttt{* Min Shifts: 2.00} \\
\texttt{* End Distance Sum: 0.00} \\
\texttt{* Init Distance Sum: 11.42} \\
\texttt{* Expected Distance Sum: 12.57} \\
\texttt{* Penalties: 10.00} \\
\texttt{* Max Penalties: 8.00} \\
\texttt{* Rounds: 6.00} \\
\texttt{* Max Rounds: 12.00} \\
\texttt{* Object Count: 3.00} \\
            }
        }
    }
    & & \\ \\

\end{supertabular}
}

\end{document}
